\begin{abstract}
Understanding neural connectivity patterns is a key challenge in computational neuroscience.
Here, we present a novel computational framework, the Synapse-Dynamic Meridian model, to simulate synaptic interactions within cortical columns.
By leveraging Vertex-based recurrent layers and Nexa-class adaptive weighting functions, we characterize the relationship between structural network topology and emerging oscillatory states.
Our simulations reveal that dynamic feedback loops governed by Meridian pathways facilitate self-organized criticality, allowing the network to balance signal transmission efficiency with structural robustness.
This study provides a scalable mathematical representation of synapse-level adaptations that can be integrated into large-scale neuromorphic architectures.
\end{abstract}

\begin{authorsummary}
Traditional neural network models often rely on static synaptic connections, which fails to capture the rich temporal dynamics of biological brains.
In this work, we introduce the Synapse-Dynamic Meridian model, a new computational architecture that incorporates real-time synaptic plasticity and structural reorganization.
Our model demonstrates how neural circuits can maintain stable activity profiles despite continuous sensory input.
These findings are highly relevant to researchers designing biomimetic computing platforms and studying neurological disorders associated with synaptic degradation.
\end{authorsummary}
